\subsubsection{Hilbert 空间}\label{subsec:hilbert-spaces}

Banach 空间具有由范数 (对应于 ``大小'' 的概念) 及其诱导的度量 (对应于 ``距离'' 的概念) 所产生的几何结构. 现在定义推广 ``角度'' 概念的内积, 以引入更多几何结构.

\begin{Definition}
\label{def:inner-product}
设 $H$ 为 $\mathbb{R}$ (或相应地, $\mathbb{C}$) 上的向量空间. \textbf{内积} $(x,y)$ 是从 $H\times H$ 到 $\mathbb{R}$ ($\mathbb{C}$) 的映射, 具有以下性质:\begin{enumerate}
\item 对每个 $x\in H$, 映射 $y\mapsto(x,y)$ 线性.
\item 对每个 $x,y\in H$, 有 $(y,x)=\overline{(x,y)}$.
\item 对每个 $x\in H$, 有 $(x,x)\geq0$, 且等号成立当且仅当 $x=0$.
\end{enumerate}
\end{Definition}

内积空间是赋范向量空间的特殊情形. 这一点由以下引理表述.

\begin{Lemma}
\label{lem:inner-product-induced-norm}
设 $H$ 为具有内积 $(\cdot,\cdot)$ 的向量空间. 那么对每个 $x,y\in H$, 有
\begin{equation}\label{eq:inner-product-cauchy-schwarz-inequality}
|(x,y)|^2\leq(x,x)(y,y).
\end{equation}
此外, 表达式
\begin{equation}\label{eq:inner-product-induced-norm}
\|x\|=\sqrt{(x,x)}
\end{equation}
定义了 $H$ 上的范数.
\end{Lemma}

不等式 \eqref{eq:inner-product-cauchy-schwarz-inequality} 称为 Cauchy-Schwarz 不等式.

\begin{Proof}
若 $x$ 或 $y$ 为 0, 则 \eqref{eq:inner-product-cauchy-schwarz-inequality} 显然成立; 现假设 $y\ne0$. 注意
\begin{equation}\label{eq:cauchy-schwarz-positive-quadratic-form}
(x+\lambda y,x+\lambda y)=(x,x)+\overline\lambda(y,x)+\lambda(x,y)+|\lambda|^2(y,y)\geq0.
\end{equation}
令 $\lambda=-(y,x)/(y,y)$ 即得 Cauchy-Schwarz 不等式. 利用 Cauchy-Schwarz 不等式, 得到
\begin{equation}\label{eq:inner-product-norm-triangle-calculation}
\begin{aligned}
\|x+y\|^2&=(x+y,x+y)\\
&=(x,x)+(x,y)+(y,x)+(y,y)\\
&\leq\|x\|^2+2\|x\|\|y\|+\|y\|^2\\
&=(\|x\|+\|y\|)^2.
\end{aligned}
\end{equation}
因此三角不等式成立. 范数的其他性质显然成立.
\end{Proof}

\begin{Definition}
\label{def:hilbert-space}
\textbf{Hilbert 空间}是 (作为赋范向量空间) 完备的内积空间.
\end{Definition}
